\documentclass[10pt,a4paper]{amsart}
\usepackage[margin=19mm]{geometry}
\usepackage{amsmath,amssymb,mathtools}
\usepackage[scr=rsfs,cal=euler]{mathalfa}
\usepackage{xfrac}
\usepackage[unicode,hidelinks,bookmarks=false]{hyperref}
\newcommand{\ie}{i.e.\ }
\newcommand{\cf}{cf.\ }
\newcommand{\resp}{respectively\ }
\newcommand{\Subsection}{\subsection}
\providecommand{\nobreakdash}{\nobreak\hbox{-}}
\setlength{\emergencystretch}{2em}
\multlinegap0pt
\usepackage{amssymb}
\newtheorem{theorem}{Theorem}
\newcommand{\Mod}[1]{\ (\mathrm{mod}\ #1)}
\title{Self-dual double cyclic codes over $\mathbb{F}_q$}
\author{Ricky Aditya, Aleams Barra, Djoko Suprijanto}
\date{Source: arXiv:2607.26823v1 (CC BY 4.0)}
\begin{document}
\maketitle

\noindent\emph{Source and licence.} Self-dual double cyclic codes over $\mathbb{F}_q$. Version arXiv:2607.26823v1 (CC BY 4.0), distributed by arXiv under the Creative Commons Attribution 4.0 International licence, http://creativecommons.org/licenses/by/4.0/, as recorded in the arXiv licence metadata for this exact version and shown on the abstract page https://arxiv.org/abs/2607.26823v1. Full licence notice: \texttt{licenses/arxiv-2607.26823-CC-BY-4.0.txt}.\par\smallskip

\noindent\textit{Editorial scope.} Counted unit: the article's theorem charsddc (source0.tex lines 177--184). The article's complete original proof of it is delivered with it (source0.tex lines 185--211, one \texttt{proof} environment of 27 lines). No mathematics is written, repaired or re-derived: the statement and the proof are the article's own lines, in the article's own notation, and no retained line is substituted or re-flowed.  Retained uncounted as the source-local context the proof needs: lines 105--112.  Nothing was omitted from any retained span, so the package carries no omission record.  No retained line contains a citation, so the wrapper carries no bibliography and no bibliography entry.  No retained line contains a figure environment, an image-inclusion command or drawing code, so no image is delivered and the package declares no asset dependency.  Editorial formatting only, in addition: an AMS \texttt{amsart} preamble that restates the article's own declarations and macros used by the retained lines, namely the environments \texttt{theorem} on separate counters, together with the standard packages those lines need.  The retained environments print in this document as Theorem 1 in order of appearance, whereas the article prints the same items with the numbering of its own full document.  The complete native source of the article as distributed in the arXiv e-print for this exact version is bundled unchanged as \texttt{sources/206294/source0.tex}.
	\begin{theorem}\label{gendual}
		If $C=\langle(b(x),0),(l(x),a(x))\rangle$ is a double cyclic code over $\mathbb{F}_q$ and $C^{\perp}=\langle(\bar{b}(x),0),(\bar{l}(x),\bar{a}(x))\rangle$ is its dual code, then
		\begin{align*}
			\bar{b}^{*}(x)&=\frac{x^r-1}{\gcd(b(x),l(x))}, \quad\bar{l}^{*}(x)=\gamma(x)\frac{x^r-1}{b(x)}, \\
			\text{ and }&\quad\bar{a}^{*}(x)=\frac{\gcd(b(x),l(x))(x^s-1)}{b(x) a(x)},&
		\end{align*}
		for some $\gamma(x)\in\mathbb{F}_q[x]$.
	\end{theorem}

	\begin{theorem}\label{charsddc}
		Let $C$ be a double cyclic code of length $(r,s)$ over $\mathbb{F}_q$ generated by $(b(x),0)$ and $(l(x),a(x))$, and let $b(x)=d(x)\beta(x)$ and $l(x)=d(x)\lambda(x)$, where $d(x)=\gcd(b(x),l(x))$. The code $C$ is self-dual if and only if these three conditions are satisfied:
		\begin{itemize}
			\item $d(x) b^{*}(x) = x^r-1$,
			\item $a(x)a^{*}(x)\beta(x)=x^s-1$,
			\item $\beta(x)$ divides $1-\lambda(x)\lambda^{*}(x)$.
		\end{itemize}
	\end{theorem}

	\begin{proof}
		($\Rightarrow$) Let the dual code $C^{\perp}$ be generated by $(\bar{b}(x),0)$ and $(\bar{l}(x),\bar{a}(x))$. If $C$ is self-dual, then $C$ and $C^{\perp}$ have the same generating sets. Using Theorem \ref{gendual}, from $b(x)=\bar{b}(x)$ we have $$b^{*}(x)=\frac{x^r-1}{\gcd(b(x),l(x))}\Leftrightarrow \gcd(b(x),l(x)) b^{*}(x) = d(x) b^{*}(x)= x^r-1,$$ and from $a(x)=\bar{a}(x)$ we have $$
		a^{*}(x)=\frac{\gcd(b(x),l(x))(x^s-1)}{b(x) a(x)} \Leftrightarrow \frac{a(x)a^{*}(x)b(x)}{\gcd(b(x),l(x))}=a(x)a^{*}(x)\beta(x)=x^s-1.$$
		Thus, the first two conditions are satisfied. Moreover, from $l(x)=\bar{l}(x)$ we have $$l^{*}(x)=\gamma(x) \frac{x^r-1}{b(x)}\Leftrightarrow b(x)l^{*}(x)=\gamma(x) (x^r-1).$$ Since $d(x)b^{*}(x)=d(x)d^{*}(x)\beta^{*}(x)=x^r-1$, taking the reciprocal of both sides, we have $d(x)d^{*}(x)\beta(x)=-(x^r-1)$. As a consequence, $$d(x)\beta(x) d^{*}(x)\lambda^{*}(x)=-\gamma(x)d^{*}(x)d(x)\beta(x)\Leftrightarrow \lambda^{*}(x)=-\gamma(x).$$
		Recall that $\gamma(x)$ is a polynomial over $\mathbb{F}_q$ such that $b(x)$ divides $d(x)+\gamma(x)l(x)$. This is equivalent to $\beta(x)$ divides $1+\gamma(x)\lambda(x)=1-\lambda(x)\lambda^{*}(x)$, i.e. the third condition is satisfied.
		
		($\Leftarrow$) First, consider the case where $\beta(x)$ is a nonzero constant polynomial and therefore a unit in $\mathbb{F}_q[x]$. In this case, $b(x)$ divides $d(x)$, and consequently also divides $l(x)=d(x)\lambda(x)$. Since $-\beta(x)^{-1}\lambda(x)(b(x),0)+(l(x),a(x))=(0,a(x))$, it follows that the code generated by $(b(x),0)$ and $(l(x),a(x))$ is equal to the code generated by $(b(x),0)$ and $(0,a(x))$. Thus, it is a separable double cyclic code. From $d(x)b^*(x)=d(x)d^*(x)\beta^*(x)=x^r-1$ and $a(x)a^{*}(x)\beta(x)=x^s-1$, we have $b(x)b^{*}(x)\Mod{x^r-1}=0$ and $a(x)a^{*}(x)\Mod{x^s-1}=0$. Thus, $(b(x),0)$ and $(0,a(x))$ form an orthogonal set and hence the code is self-dual.
		
		Now consider the other case where $\beta(x)$ is not a constant polynomial. Let $m=\operatorname{lcm}(r,s)$. Recall that $\{(b(x),0),(l(x),a(x))\}$ is an orthogonal set if and only if these three conditions are satisfied:
		\begin{itemize}
			\item $b(x)b^{*}(x)\Mod{x^r-1}=0$,
			\item $b(x)l^{*}(x)\Mod{x^r-1}=0$,
			\item $\left[l(x)l^{*}(x)\Mod{x^r-1}\right] \left(\frac{x^m-1}{x^r-1}\right)+\left[a(x)a^{*}(x)\Mod{x^s-1}\right]\left(\frac{x^m-1}{x^s-1}\right)=0$.
		\end{itemize}
		Since $d(x) b^{*}(x) = x^r-1$, we have $b(x)b^{*}(x)=\beta(x) d(x)b^{*}(x)=\beta(x)(x^r-1)$ and $b(x)l^{*}(x)=b(x) d^{*}(x)\lambda^{*}(x)=-\lambda^{*}(x)(x^r-1)$. In other words, $b(x)b^{*}(x)$ and $b(x)l^{*}(x)$ are both divided by $x^r-1$. These satisfy the first two conditions.
		
		From $a(x)a^{*}(x)\beta(x)=x^s-1$ and $\beta(x)$ being a nonconstant polynomial, it is clear that $\deg(a(x)a^{*}(x))<s$. Since $\beta(x)$ divides $1-\lambda(x)\lambda^{*}(x)$, we can write $1-\lambda(x)\lambda^{*}(x)=\beta(x)\nu(x)$, for some $\nu(x)\in\mathbb{F}_q[x]$. Then we obtain
		\begin{align*}
			\left[l(x)l^{*}(x)\Mod{x^r-1}\right]&\left(\frac{x^m-1}{x^r-1}\right)+\left[a(x)a^{*}(x)\Mod{x^s-1}\right]\left(\frac{x^m-1}{x^s-1}\right)\\
			&=\left[l(x)l^{*}(x)\Mod{x^r-1}\right]\left(\frac{x^m-1}{x^r-1}\right)+\frac{x^s-1}{\beta(x)}\frac{x^m-1}{x^s-1}\\
			&=\left[\left(l(x)l^{*}(x)+\frac{x^r-1}{\beta(x)}\right)\Mod{x^r-1}\right]\left(\frac{x^m-1}{x^r-1}\right)\\
			&=\left[d(x)d^{*}(x)(\lambda(x)\lambda^{*}(x)-1)\Mod{x^r-1}\right]\left(\frac{x^m-1}{x^r-1}\right)\\
			&=\left[-d(x)d^{*}(x)\beta(x)\nu(x)\Mod{x^r-1}\right]\left(\frac{x^m-1}{x^r-1}\right)\\
			&=\left[\nu(x)(x^r-1)\Mod{x^r-1}\right]\left(\frac{x^m-1}{x^r-1}\right)=0
		\end{align*}
		This satisfies the third condition. Therefore, $C=\langle (b(x),0),(l(x),a(x))\rangle$ is a self-dual double cyclic code of length $(r,s)$ over $\mathbb{F}_q$.
	\end{proof}

\end{document}
